% lsqvce -- the mathematical framework. Built with pdflatex (build.sh in this folder).
% Created 2026-10-12 (CEST, UTC+02:00) with assistance from Claude (Anthropic),
% model Claude Opus 5.5 -- WP-124. The September 2026 R5 document it replaces is
% archived under dev/reviews/R5_2026-09-01/theory_r5/.
\documentclass[11pt,a4paper,twoside,openany]{report}
\usepackage[T1]{fontenc}
\usepackage[utf8]{inputenc}
\usepackage{mathptmx}\usepackage[scaled=0.92]{helvet}\usepackage{courier}
\usepackage{microtype}
\usepackage[margin=25mm,headheight=14pt]{geometry}
\usepackage{amsmath,amssymb,bm}
\usepackage{booktabs,array}
\usepackage{xcolor}
\usepackage{enumitem}
\usepackage[most]{tcolorbox}
\usepackage{titlesec}
\usepackage{fancyhdr}
\usepackage[hidelinks,pdftitle={lsqvce -- Mathematical framework},pdfauthor={Matthias Weigelt}]{hyperref}

\setlist{nosep,leftmargin=*}
\sloppy
\setlength{\parskip}{4pt}
\setlength{\parindent}{0pt}
\titleformat{\chapter}[hang]{\bfseries\LARGE}{\thechapter}{1em}{}
\titlespacing*{\chapter}{0pt}{-10pt}{14pt}
\numberwithin{equation}{chapter}

\definecolor{boxbg}{RGB}{245,245,240}
\definecolor{rule}{RGB}{120,120,120}
\newtcolorbox{remark}[1][]{enhanced,colback=boxbg,colframe=rule,boxrule=0.5pt,arc=1pt,left=6pt,right=6pt,top=4pt,bottom=4pt,fonttitle=\bfseries,title={#1}}

\newcommand{\mat}[1]{\bm{#1}}
\newcommand{\vc}[1]{\bm{#1}}
\newcommand{\T}{^{\mathsf T}}
\newcommand{\tr}{\operatorname{tr}}
\newcommand{\diag}{\operatorname{diag}}
\newcommand{\E}{\operatorname{E}}
\newcommand{\Var}{\operatorname{Var}}
\newcommand{\MAD}{\operatorname{MAD}}
\newcommand{\code}[1]{\texttt{#1}}

\pagestyle{fancy}
\fancyhf{}
\fancyhead[LE,RO]{\small\thepage}
\fancyhead[RE]{\small\nouppercase{\leftmark}}
\fancyhead[LO]{\small lsqvce -- mathematical framework}
\renewcommand{\headrulewidth}{0.3pt}

\begin{document}

\begin{titlepage}
\vspace*{35mm}
{\bfseries\Huge lsqvce\par}
\vspace{6mm}
{\bfseries\LARGE Least-squares adjustment with variance component estimation\par}
\vspace{4mm}
{\Large The mathematical framework\par}
\vspace{20mm}
\begin{remark}[Scope]
This document states what the toolbox computes and why: the functional model and its
solution, the stochastic model, the estimation of variance components, robust and
nonlinear estimation, and the quality measures. How to use the toolbox is the subject
of the Manual (\code{doc lsqvce}); the interfaces are specified in the API reference.
Equation numbers of the form ``Ellmer (6.5.x)'' refer to Ellmer (2018).
\end{remark}
\vfill
{\small Matthias Weigelt, with assistance from Claude (Anthropic).\par}
\end{titlepage}

\tableofcontents

% Chapter 1 -- the functional model and its solution.
\chapter{The linear model and its solution}\label{ch:model}

\section{Gauss--Markov model}

The observations $\vc l\in\mathbb R^n$ are modelled as
\begin{equation}\label{eq:gm}
\vc l = \mat A\vc x + \vc e,\qquad \E\{\vc e\}=\vc 0,\qquad \mathrm{C}\{\vc e\}=\sigma_0^2\mat\Sigma ,
\end{equation}
with the design matrix $\mat A$ ($n\times u$, full column rank unless stated otherwise), the
unknown parameters $\vc x$, the cofactor matrix $\mat\Sigma$ and the variance of unit weight
$\sigma_0^2$. The best linear unbiased estimate and its covariance are
\begin{equation}\label{eq:lsq}
\hat{\vc x}=\mat N^{-1}\vc n,\quad \mat N=\mat A\T\mat P\mat A,\quad \vc n=\mat A\T\mat P\vc l,\quad
\mat P=\mat\Sigma^{-1},\qquad
\hat\sigma_0^2=\frac{\hat{\vc v}\T\mat P\hat{\vc v}}{n-u},\quad
\mat\Sigma_{\hat x\hat x}=\hat\sigma_0^2\mat N^{-1},
\end{equation}
with the residuals $\hat{\vc v}=\vc l-\mat A\hat{\vc x}$. The residual is a projection,
$\hat{\vc v}=(\mat I-\mat H)\vc l$ with the hat matrix $\mat H=\mat A\mat N^{-1}\mat A\T\mat P$, and
the redundancy $n-u=\tr(\mat I-\mat H)$ is distributed over the observations by the diagonal of
$\mat R=\mat I-\mat H$, $\sum_i r_i=n-u$. Every variance estimator of Chapter~\ref{ch:vce} is a
specialisation of ``quadratic form over its expectation''.

\section{Whitening}\label{sec:whiten}

With a factor $\mat W$ of $\mat P=\mat W\T\mat W$ the model becomes
$\bar{\vc l}=\bar{\mat A}\vc x+\bar{\vc e}$ with $\bar{\mat A}=\mat W\mat A$, $\bar{\vc l}=\mat W\vc l$ and
$\mathrm C\{\bar{\vc e}\}=\sigma_0^2\mat I$, and \eqref{eq:lsq} is ordinary least squares. The
toolbox takes $\mat W=\mat L^{-1}$, the inverse of the lower Cholesky factor of
$\mat\Sigma=\mat L\mat L\T$. $\mat\Sigma$ is never formed for all observations: observations are
grouped into \emph{arcs}, independent of each other, and every arc is whitened by its own law
(Chapter~\ref{ch:stochastic}). All quantities below --- normal equations, residuals,
redundancies, test statistics --- are expressed in the whitened frame with a-priori
$\sigma_0=1$; ``residual'' means $\bar{\vc v}=\mat W\hat{\vc v}$ unless said otherwise.

\section{Two ways to accumulate and solve}\label{sec:backends}

\paragraph{Square-root information (QR, the default).} The upper-triangular factor $\mat R$ of
$[\bar{\mat A}\ \bar{\vc l}]$ is updated batch by batch by Householder reduction,
$\mat R\leftarrow\operatorname{triu}\operatorname{qr}\bigl([\mat R;\ \bar{\mat A}_k\ \bar{\vc l}_k]\bigr)$, and
truncated to $u+1$ rows. With
\[
\mat R=\begin{bmatrix}\mat R_1 & \vc r\\ \vc 0\T & \rho\end{bmatrix},\qquad
\hat{\vc x}=\mat R_1^{-1}\vc r,\qquad \bar{\vc v}\T\bar{\vc v}=\rho^2,\qquad
\mat\Sigma_{\hat x\hat x}=\hat\sigma_0^2(\mat R_1\T\mat R_1)^{-1},
\]
so the residual sum of squares is obtained without subtraction, and the condition number is
that of $\bar{\mat A}$, not its square. The rank decision is a pivot test: the system is
refused when $\min_k|R_{1,kk}|\le 10^3\,\varepsilon\,\max_k|R_{1,kk}|$.

\paragraph{Normal equations.} $\mat N=\sum_k\bar{\mat A}_k\T\bar{\mat A}_k$,
$\vc n=\sum_k\bar{\mat A}_k\T\bar{\vc l}_k$ and $\bar{\vc l}\T\bar{\vc l}$ are summed, and
$\mat N=\mat U\T\mat U$ is factorised by Cholesky. Partial systems add, which is the exchange
format of the field, but $\kappa(\mat N)=\kappa(\bar{\mat A})^2$ and
\begin{equation}\label{eq:cancel}
\bar{\vc v}\T\bar{\vc v}=\bar{\vc l}\T\bar{\vc l}-\vc n\T\hat{\vc x}
\end{equation}
cancels catastrophically when the fit is very good relative to $\|\bar{\vc l}\|$. Where the
cancellation exceeds the attainable precision the value is not trusted: the residual sum is
recomputed from the residuals in a second pass over the data.

\section{Block-arrow systems and the elimination of local parameters}\label{sec:blockarrow}

The parameters split into \emph{common} parameters $\vc x_c$ ($n_c$), seen by every observation,
and \emph{local} parameters $\vc y_s$ ($n_{l,s}$) of system $s$ --- initial states, instrument
calibrations, empirical terms --- valid for one arc, one day or the whole system. For the rows of
system $s$
\begin{equation}\label{eq:blockarrow}
\bar{\vc l}_s=\bar{\mat A}_s\vc x_c+\bar{\mat B}_s\vc y_s+\bar{\vc e}_s .
\end{equation}
The normal matrix has block-arrow form. The local block of each system --- all its arc-, day-
and system-scoped parameters together --- is eliminated by the Schur complement
\begin{equation}\label{eq:schur}
\mat N_{cc}^{\mathrm{red}}=\mat N_{cc}-\mat N_{lc}\T\mat N_{ll}^{-1}\mat N_{lc},\qquad
\vc n_c^{\mathrm{red}}=\vc n_c-\mat N_{lc}\T\mat N_{ll}^{-1}\vc n_l,\qquad
(\bar{\vc l}\T\bar{\vc l})^{\mathrm{red}}=\bar{\vc l}\T\bar{\vc l}-\vc n_l\T\mat N_{ll}^{-1}\vc n_l ,
\end{equation}
and the reduced systems of all systems add. On the square-root backend the same reduction is a
re-triangularisation with the local columns first,
\[
\begin{bmatrix}\mat R_{11}&\mat R_{12}&\vc r_1\\ & \mat R_{22}&\vc r_2\\ && \rho\end{bmatrix},
\]
whose trailing rows $[\mat R_{22}\ \vc r_2;\ \vc 0\T\ \rho]$ are the reduced contribution, and
$\mat R_{11},\mat R_{12},\vc r_1$ give the local solution and cofactor without forming
$\mat N_{ll}$. Each system contributes $m_s-n_{l,s}$ degrees of freedom. A local block that fails
the pivot test is rank-deficient: with $\mat R_{11}=\mat U\mat S\mat V\T$ and $k$ the singular values
above the rank tolerance, the locals live on $\mat V_1$ (the minimum-norm $\hat{\vc y}$, with
$\mat N_{ll}^{-1}$ read as the pseudo-inverse $\mat V_1\mat S_1^{-2}\mat V_1\T$ below), the rows
$\mat U_2\T[\mat R_{12}\ \vc r_1]$ join the reduced contribution, and the system contributes
$m_s-k$ degrees of freedom. Forming $\mat N_{ll}$ and cutting its eigenvalues instead loses a
weak but determined direction together with the null one.

Back-substitution recovers
\begin{equation}\label{eq:backsub}
\hat{\vc y}_s=\mat N_{ll}^{-1}(\vc n_l-\mat N_{lc}\hat{\vc x}_c),\qquad
\mat\Sigma_{\hat y\hat y}=\hat\sigma_0^2\mat N_{ll}^{-1}+\mat G\mat\Sigma_{\hat x\hat x}\mat G\T,\quad
\mat G=\mat N_{ll}^{-1}\mat N_{lc}.
\end{equation}
Every diagnostic that involves the hat matrix includes the local block. With
$\vc a=[\vc a_c;\vc b]$ a row in the full parameter space and $\vc z=\vc a_c-\mat G\T\vc b$,
\begin{equation}\label{eq:leverage}
\vc a\T\mat Q_{xx}\vc a=\bigl\|\mat R_c^{-\mathsf T}\vc z\bigr\|^2+\bigl\|\mat H_s\vc b\bigr\|^2,\qquad
\mat H_s\T\mat H_s=\mat N_{ll}^{-1},
\end{equation}
where $\mat R_c$ is the factor of the reduced common system; the full $u\times u$ cofactor matrix
is never formed.

\section{Prior information and regularisation}\label{sec:priors}

A prior is a group $k$ of $p_k$ pseudo-observations
\begin{equation}\label{eq:prior}
\mat L_k\vc x=\vc x_{0,k}+\vc e_k,\qquad \mathrm C\{\vc e_k\}=\sigma_k^2\mat I ,
\end{equation}
written in the columns of the parameters. In the whitened frame it contributes the rows
$\mat L_k/\sigma_k$ with observations $\vc x_{0,k}/\sigma_k$; this is Tikhonov regularisation with
$\lambda_k=\sigma_0^2/\sigma_k^2$, $\mat N_\lambda=\mat N+\sum_k\mat L_k\T\mat L_k/\sigma_k^2$.
Relative weights inside a group are carried by the rows of $\mat L_k$; $\sigma_k$ is one common
factor and can be estimated (\S\ref{sec:priorvce}).

Priors on common parameters are stacked into one operator $\mat K=[\mat L_k/\sigma_k]_k$ and
applied at the solve: on the square-root backend by appending its rows to the factor, on the
normal backend as $\mat N+\mat K\T\mat K$. A prior on local parameters is an observation stream of
its own with $\bar{\mat A}=\vc 0$, $\bar{\mat B}=\mat L/\sigma$ and $\bar{\vc l}=\vc x_0/\sigma$, so it
enters the elimination \eqref{eq:schur} like data. Either way the prior rows count as
observations: the redundancy is $n+\sum_k p_k-u$, and the share the data overrule is the
redundancy of the group,
\begin{equation}\label{eq:priorred}
r_k=p_k-\tr\bigl(\mat L_k\mat N_\lambda^{-1}\mat L_k\T\bigr)/\sigma_k^2 .
\end{equation}

\subsection{The datum of a free network}\label{sec:datum}
A rank defect --- a datum that no observation fixes --- is removed by a prior on the null space.
Let $\mat K_g$ be the rows of group $g$. The data are \emph{blind} to the group when
\[
\|\mat N\mat K_g\T\|_F\le\epsilon\,\bigl\|\,|\mat N|\,|\mat K_g\T|\,\bigr\|_F
\]
(elementwise absolute values: the product against its own rounding bound, which does not change
with the units of the parameters; on the square-root backend $\mat N$ is replaced by the factor).
A blind group is a \emph{datum} when it is also orthogonal to the rows $\mat K_o$ of every group
the data see, $\|\mat K_o\mat K_g\T\|_F\le\epsilon\,\|\mat K_o\|_F\|\mat K_g\|_F$; a datum may be
split over several groups. With $\mat K_d$ the rows
of all datum groups and $\mat K_i$ the rest, $\mat N_i=\mat N+\mat K_i\T\mat K_i$. Because the range
of the data and the null space are orthogonal, the estimate $\hat{\vc x}$ is the
minimum-constraint solution for any datum weight, and its covariance is that of the
minimum-constraint problem,
\begin{equation}\label{eq:mincon}
\mat\Sigma_{\hat x\hat x}=\hat\sigma_0^2\,\mat N_i^{+},
\end{equation}
evaluated stably as $\hat\sigma_0^2(\mat R_m^{-1}\mat R_m^{-\mathsf T}-\mat D_m\mat D_m\T)$ with
$\mat R_m\T\mat R_m=\mat N_i+c^2\mat K_d\T\mat K_d$, $\mat D_m=\mat K_d\T(\mat K_d\mat K_d\T)^{-1}/c$ and a
balancing constant $c$. A group only partly in the null space, or one that overlaps another group,
is information: its estimate moves with its $\sigma$, and it keeps the covariance
$\hat\sigma_0^2(\mat N+\mat K\T\mat K)^{-1}$.

\section{Nonlinear models: damped Gauss--Newton}\label{sec:gn}

For $\vc l=f(\vc x)+\vc e$ the objective, prior included, is
\begin{equation}\label{eq:gnobj}
F(\vc x)=\bigl\|\mat W(\vc l-f(\vc x))\bigr\|^2+\sum_k\bigl\|(\mat L_k\vc x-\vc x_{0,k})/\sigma_k\bigr\|^2 .
\end{equation}
At the iterate $\vc x$ the Jacobian $\mat J$ is accumulated once; that accumulation is both the
linearisation and the evaluation of $F(\vc x)$. The step $\Delta\vc x$ is one least-squares solve
with two kinds of pseudo-rows on the common block:
\begin{equation}\label{eq:gnstep}
\mat K\,\Delta\vc x\approx\vc x_0-\mat K\vc x,\qquad
\sqrt{\lambda d_j}\,\Delta x_j\approx0,\quad \vc d=\diag(\mat J\T\mat P\mat J).
\end{equation}
The first form makes the iteration converge to the minimiser of \eqref{eq:gnobj}; the second is
Marquardt damping of the parameters themselves. A step is accepted only if
$F(\vc x+\Delta\vc x)\le F(\vc x)$; a rejected step multiplies $\lambda$ by ten (from zero to
$10^{-3}$) and is re-solved on the same accumulation, an accepted one divides it by ten (below
$10^{-12}$ it is set to zero). The iteration stops when $\max_j|\Delta x_j|/\sigma_{x_j}$ or the
relative change of $F$ falls below its tolerance. The covariance is that of the undamped
linearised problem at the final point, so it does not depend on the damping history; a run that
exhausts its budget returns the best accepted point.

\section{What the estimates mean}\label{sec:meaning}

$\mat\Sigma_{\hat x\hat x}=\hat\sigma_0^2\mat N^{-1}$ is the \emph{formal} covariance: correct if the
functional model is complete, the stochastic model is right and the residuals are white.
Chapter~\ref{ch:vce} makes the second condition true as far as the data allow;
Chapter~\ref{ch:quality} tests the first and the third.

% Chapter 2 -- the stochastic model.
\chapter{The stochastic model}\label{ch:stochastic}

\section{Stationary noise per arc}

Within one arc of $m$ equally spaced epochs the noise of a scalar stream is modelled as a
stationary process with autocovariance $c_\tau$. Its covariance matrix is Toeplitz,
\begin{equation}\label{eq:toeplitz}
\mat\Sigma=\bigl[c_{|i-j|}\bigr]_{i,j=1}^m ,
\end{equation}
and it is positive definite for every $m$ if and only if the spectrum of the process is positive
(Bochner). Arcs are independent of each other; non-stationarity across a campaign is carried by
one variance factor per arc (\S\ref{sec:arcfactors}).

\section{Temporal laws}

A law is declared by one of the following, each with a level $\sigma$ (a standard deviation):

\begin{center}
\begin{tabular}{@{}ll@{}}
\toprule
law & covariance \\ \midrule
white & $\mat\Sigma=\diag(\sigma_i^2)$ \\
ARMA$(p,q)$ & $c_\tau$ of the process $\sum_{k=0}^p a_ke_{t-k}=\sum_{k=0}^q b_k\eta_{t-k}$, $\Var\eta=\sigma^2$\\
covariance function & $\Sigma_{ij}=\sigma^2\rho(t_i-t_j)$ (actual time differences) \\
spectrum & $c_\tau=\int_0^{f_{\mathrm N}}S(f)\cos(2\pi f\tau\Delta t)\,\mathrm df$ \\
vector AR$(p)$ & $\vc e_t=\sum_{k=1}^p\mat A_k\vc e_{t-k}+\vc\eta_t$, $\mathrm C\{\vc\eta\}=\mat\Sigma_\eta$ \\
\bottomrule
\end{tabular}
\end{center}

\paragraph{Spectrum to covariance.} A spectrum is given by ordinates $S_j$ at frequencies $f_j$
and is taken as piecewise linear between them. Its autocovariance is the exact integral; for one
segment $[a,b]$ with centre $c$, half-width $h$ and $x=2\pi\tau h$,
\begin{equation}\label{eq:psd2acv}
\int_a^b S(f)\cos(2\pi f\tau)\,\mathrm df=(S_a+S_b)\,h\cos(2\pi\tau c)\frac{\sin x}{x}
-(S_b-S_a)\,h\sin(2\pi\tau c)\frac{\sin x-x\cos x}{x^2}.
\end{equation}
The result is valid at every lag; a non-negative spectrum gives a positive semi-definite
covariance of any size.

\section{Whitening an arc}\label{sec:whitening}

\paragraph{Dense factor.} $\mat W=\mat L^{-1}$ with $\mat L\mat L\T=\mat\Sigma$, cost $O(m^3)$ and
memory $O(m^2)$ per arc. Exact for any positive definite law. The factor is formed over the whole
arc: blocks of a dense factor cannot be joined by an overlap, because the conditional distribution
given a finite past is not the conditional given the whole past.

\paragraph{Autoregressive filter.} For an AR$(p)$ process with coefficients
$\vc a=[1,a_1,\dots,a_p]$ and innovation standard deviation $s$, whitening is the FIR filter
\begin{equation}\label{eq:arfilter}
\bar x_t=\frac1s\sum_{k=0}^p a_kx_{t-k},\qquad t>p,
\end{equation}
cost $O(mp)$, with no covariance matrix formed. The first $p$ rows are replaced by the exact
stationary start $\mat L_p^{-1}\vc x_{1:p}$, $\mat L_p\mat L_p\T=\operatorname{toeplitz}(c_0,\dots,c_{p-1})$.
The resulting operator is the inverse Cholesky factor of $\mat\Sigma$ itself: both are lower
triangular square roots of $\mat\Sigma^{-1}$, which is unique up to sign. Consecutive batches
overlap by the filter memory, so batching does not change the result. For an ARMA$(p,q)$ law the
whitener $A(z)/B(z)$ is recursive; its memory is $K=p+N$, with $N$ the number of samples after
which the impulse response of $1/B(z)$ has decayed below the rounding level,
$|h_N|\le\varepsilon\max|h|$ (finite because $B(z)$ is required to have all roots inside the
unit circle). The recursion carries state, so the exact start covers the first $K$ rows,
$\mat L_K^{-1}\vc x_{1:K}$, and continues from their exact innovations
$u_k=L_{kk}\bar x_k$:
\begin{equation}
u_k=\sum_{j=0}^p a_jx_{k-j}-\sum_{j=1}^q b_ju_{k-j},\qquad \bar x_k=u_k/s,\qquad k>K;
\end{equation}
beyond $K$ the coefficients of the innovations algorithm equal $\vc b$ to rounding, so the
operator is the inverse Cholesky factor of $\mat\Sigma$. The AR coefficients of a
law given by its autocovariance follow from the Levinson--Durbin recursion (Yule--Walker); a
high-order AR can be reduced to an ARMA. A spectrum with narrow features is represented better by
the dense factor than by a low-order filter.

\section{Multi-axis streams and the per-epoch accuracy}\label{sec:multiaxis}

A stream with $d$ axes per epoch is ordered epoch-major. What the preprocessing knows about the
accuracy of epoch $e$ --- per-axis standard deviations or a full $d\times d$ covariance
$\mat C_e$ --- enters as the factor $\mat D_e$ ($\mat D_e=\diag(\sigma_{e,1},\dots)$, or the lower
Cholesky factor of $\mat C_e$). With the temporal law $\mat R_t$ on the epoch grid, the covariance of
an arc is
\begin{equation}\label{eq:separable}
\mat\Sigma=\mat D\,(\mat R_t\otimes\mat I_d)\,\mat D\T,\qquad \mat D=\operatorname{blkdiag}(\mat D_e),
\end{equation}
and its whitener factorises: every epoch is first de-scaled by $\mat D_e^{-1}$, then the temporal
whitener runs down the epochs of each axis. The order matters: the de-scaled errors are
stationary, the raw ones are not when $\mat D_e$ varies. Without per-epoch information
$\mat D=\mat I$. For a scalar stream ($d=1$) the same composition applies with $\mat D_e=\sigma_e$.

Cross-axis correlation \emph{in time} is outside \eqref{eq:separable}; it is the vector AR model,
fitted by the multichannel Levinson (Whittle) recursion and whitened by the matrix filter
$\vc w_t=\mat L_\eta^{-1}\bigl(\vc z_t-\sum_k\mat A_k\vc z_{t-k}\bigr)$ with an exact block start.

\section{Where a law can be wrong}\label{sec:nonpd}

A covariance function truncated at a lag shorter than the arc is in general indefinite: its
spectrum is the true spectrum convolved with a kernel that has negative side lobes. A law defined
by a non-negative spectrum cannot fail this way. An indefinite law is treated as a fault: in
accumulation its eigenvalues are floored with a warning; in the variance-component sweep the arc
is skipped (\S\ref{sec:skips}).

% Chapter 3 -- variance component estimation (the joint estimator).
\chapter{Variance component estimation}\label{ch:vce}

\section{The problem}

The stochastic model is known up to a finite set of non-negative components entering linearly,
\begin{equation}\label{eq:vcmodel}
\mat\Sigma(\vc\theta)=\sum_{g=1}^G\theta_g\mat V_g ,
\end{equation}
with known symmetric positive semi-definite matrices $\mat V_g$. The toolbox estimates, jointly
in one iteration,
\begin{enumerate}[label=(\alph*)]
\item the \textbf{spectral shape} of the noise of every pool group of streams
  (\S\ref{sec:hat}),
\item the \textbf{level} of every stream,
\item a \textbf{variance factor per arc} (\S\ref{sec:arcfactors}), and
\item the \textbf{weight of every prior group} declared as estimable (\S\ref{sec:priorvce}).
\end{enumerate}

\section{The Helmert estimator and its simplified form}\label{sec:helmert}

The quadratic forms of the residuals have the expectation (Koch 1999)
\begin{equation}\label{eq:helmert}
\E\bigl[\hat{\vc v}\T\mat P\mat V_g\mat P\hat{\vc v}\bigr]=\sum_h\theta_h\tr(\mat R\mat V_g\mat R\mat V_h),\qquad
\mat R=\mat P-\mat P\mat A\mat N^{-1}\mat A\T\mat P .
\end{equation}
Equating observed and expected forms gives the Helmert system, an unbiased estimator. The
simplified estimator drops the coupling between components and iterates
\begin{equation}\label{eq:simplified}
\theta_g\leftarrow\theta_g\,\frac{\Omega_g}{r_g},\qquad
\Omega_g=\theta_g\,\hat{\vc v}\T\mat P\mat V_g\mat P\hat{\vc v},\qquad
r_g=\theta_g\tr(\mat R\mat V_g).
\end{equation}
The ratio $\alpha_g^2=\Omega_g/r_g$ is the \emph{Helmert factor} of component $g$; the fixed
point $\Omega_g=r_g$ (``a weighted residual sum equals its redundancy'') is shared with the full
estimator. Every component in this chapter is estimated by \eqref{eq:simplified}; the components
differ only in $\mat V_g$.

\section{The spectral law and its components}\label{sec:hat}

\paragraph{Parameterisation.} The law of a pool group is a non-negative spectrum $\vc S$ on $L$
bins $f_j=j\,d$, $d=1/(2(L-1))$ cycles per sample, $j=0,\dots,L-1$, linear between the bins. Its
autocovariance at every lag $\tau=0,1,\dots$ of every arc is the exact integral
\begin{equation}\label{eq:hatacv}
\vc r=\mat K\vc S,\qquad
K_{\tau j}=2c_j\cos(2\pi f_j\tau)\,\frac{1-\cos(2\pi\tau d)}{d\,(2\pi\tau)^2},\qquad K_{0j}=c_jd,
\end{equation}
with $c_j=2$ inside and $c_j=1$ at both ends. Column $j$ of $\mat K$ is the autocovariance of the
``hat'' spectrum of bin $j$, so $\operatorname{toeplitz}(\mat K_{:j})$ is positive semi-definite for
every arc length, and the components
\begin{equation}
\mat V_j=\operatorname{toeplitz}(\mat K_{:j}),\qquad \theta_j=S_j\ge0,
\end{equation}
make \eqref{eq:vcmodel} a sum of positive semi-definite terms. A Helmert factor of a positive
semi-definite component cannot be negative, so the law needs no floor and no constraint. The
grid sets the spectral resolution, not the lag range: the law is exact at every lag whatever $L$.
The law is a \emph{shape}: it is normalised to $r_0=1$ after every update, and the level is
carried separately. The start is white noise, $\vc S=\vc 1$.

\paragraph{The per-lag statistics.} Write the arc covariance as
$\mat\Sigma=\lambda\sum_j S_j\mat V_j$ and let $\mat V^{(n)}$ have ones on the $\pm n$-th diagonals.
Per arc and lag the toolbox forms the residual sums and the redundancies
\begin{equation}\label{eq:lagstats}
\Omega_n=\tilde{\vc e}\T\mat V^{(n)}\tilde{\vc e},\quad\tilde{\vc e}=\mat\Sigma^{-1}\hat{\vc e},\qquad
s_n=\tr\bigl(\bar{\mat R}\mat V^{(n)}\bigr),
\end{equation}
for all lags $n=0,\dots,m-1$ of the arc, with $\bar{\mat R}$ the redundancy matrix \eqref{eq:helmert}
restricted to the arc. Since $\mat V_j=\sum_n K_{nj}\mat V^{(n)}$, the projections on the bins are
\begin{equation}\label{eq:binstats}
Q_j=\lambda\sum_nK_{nj}\Omega_n,\qquad T_j=\lambda\sum_nK_{nj}s_n ,
\end{equation}
computed for all bins at once by folding the lag sums and one FFT.

\section{The redundancies by random probes}\label{sec:hutchinson}

$\Omega_n$ is a sum of products of residuals and is exact. The redundancy needs $\mat N^{-1}$. Per
arc, with the whitened designs $\bar{\mat A}=\mat W(\mat A-\mat B\mat G)$ and $\bar{\mat B}=\mat W\mat B$, the
local block exact and the common block by Hutchinson's estimator,
\begin{equation}\label{eq:hutch}
\bar{\mat R}=\mat\Sigma^{-1}-\mat R_t\mat R_t\T-\mat K_l\mat K_l\T,\qquad
\mat R_t=\mat W\T\bar{\mat A}\,\mat U^{-1}\bar{\mat Z},\qquad
\mat K_l=\mat W\T\bar{\mat B}\,\mat L_l^{-\mathsf T},
\end{equation}
with $\mat N=\mat U\T\mat U$ the reduced common normal matrix (common priors included),
$\mat N_{ll}=\mat L_l\mat L_l\T$, and $\bar{\mat Z}=\mat Z/\sqrt Z$ a matrix of $Z$ Rademacher probe
vectors. Since $\E[\bar{\mat Z}\bar{\mat Z}\T]=\mat I$, $\E[\mat R_t\mat R_t\T]$ is the exact common
term. The $u\times Z$ block $\mat U^{-1}\bar{\mat Z}$ is formed once per pass and shared by all arcs
(Ellmer 6.5.48); the same probes are used in every pass. Then
\begin{equation}
s_0=\tr\bar{\mat R},\qquad s_n=2\sum_i\bar R_{i+n,i},
\end{equation}
where the sub-diagonal sums of $\mat R_t\mat R_t\T$ and $\mat K_l\mat K_l\T$ are evaluated for all lags
by one zero-padded FFT per column, and those of $\mat\Sigma^{-1}$ by the Levinson recursion and the
Gohberg--Semencul representation when $\mat\Sigma$ is Toeplitz (exact columns otherwise). The
estimate is unbiased; its spread decreases as $1/\sqrt Z$. There is no exact route in the joint
estimator; per-observation redundancies for the diagnostics are computed exactly by a different
route (\S\ref{sec:redundancy}).

\section{One pass}\label{sec:pass}

Each pass of the joint estimator
\begin{enumerate}
\item accumulates and solves under the current model;
\item sweeps every arc once, whitening it with its current covariance and forming
  \eqref{eq:lagstats}--\eqref{eq:binstats} and the arc's Helmert factor (\S\ref{sec:arcfactors});
\item updates the spectrum of every pool group (streams with an estimated law only),
  \begin{equation}\label{eq:psdupdate}
  \alpha_j^2=\frac{\max(\sum_mQ_{jm},0)}{\sum_mT_{jm}},\qquad S_j\leftarrow\alpha_j^2S_j ,
  \end{equation}
  summed over the arcs $m$ of the group. A bin whose redundancy $S_j\sum_mT_{jm}$ falls below a
  threshold keeps its value ($\alpha_j^2=1$); an optional moving average smooths $\vc S$; a
  relative numerical floor of $10^{-10}\max\vc S$ protects the factorisation; then $\vc S$ is
  renormalised to $r_0=1$;
\item updates the arc factors and stream levels (\S\ref{sec:arcfactors}) and the prior
  weights (\S\ref{sec:priorvce});
\item records the whiteness of every stream (\S\ref{sec:whiteness}).
\end{enumerate}
The iteration stops when
\begin{equation}
\max_j|\alpha_j^2-1|\le\varepsilon\qquad\text{and}\qquad\max_m\Bigl|\frac{\sigma^2_{m,\mathrm{new}}}{\sigma^2_{m,\mathrm{old}}}-1\Bigr|\le\varepsilon ,
\end{equation}
or after the pass budget. After the loop the system is accumulated once more under the final
model; that solution is the result.

\paragraph{Pooling.} Streams of a pool group share one spectrum $\vc S$; every stream keeps its own
level. Pooling more streams gives each bin more redundancy per pass; pooling streams of different
physics forces one shape on different spectra.

\paragraph{Policies per stream.} A stream may have its law estimated (shape, level and arc
factors), only its level and arc factors (the declared law is kept as the shape), or nothing. An
estimated law replaces the declared one from the first pass; the declared level is kept as the
multiplier of the estimated components.

\section{Arc factors and stream levels}\label{sec:arcfactors}

The covariance of arc $m$ of stream $s$ is
\begin{equation}\label{eq:arccov}
\mat\Sigma_m=\sigma_{\mathrm{decl},s}^2\;\sigma_m^2\;\operatorname{toeplitz}(\mat K\vc S_g),
\end{equation}
composed with the arc's per-epoch accuracy \eqref{eq:separable}. The Helmert factor of the arc's
whole covariance (Ellmer 6.5.27) is
\begin{equation}
\alpha_m^2=\frac{q_m}{t_m},\qquad q_m=\vc r_m\T\vc\Omega^{(m)},\quad t_m=\vc r_m\T\vc s^{(m)},
\end{equation}
with $\vc r_m$ the first column of the arc's temporal covariance, and the updated factor is
$\gamma_m^2=\alpha_m^2\sigma_m^2$ (Ellmer 6.5.28). The factors are \emph{absolute}: they weight the
solve, so $\sigma_m^2$ is the arc's variance relative to its declaration and $\hat\sigma_0$ settles
near one. A factor is clamped below at $10^{-8}\bar\gamma^2$, with $\bar\gamma^2$ the mean of the
central half of the stream's factors, so that no arc is silenced without a test having rejected it.
A stream without arc factors carries one level, the mean of its arcs' Helmert factors.

The standard deviation of an arc is therefore the product
$\sigma_{\mathrm{decl},s}\cdot\sigma_m\cdot\hat\sigma_0$.

\section{Weights of prior groups}\label{sec:priorvce}

The weight of a common prior group $k$ (\S\ref{sec:priors}) is a variance component of its own
(Koch and Kusche 2002). From the solution of the previous pass,
\begin{equation}\label{eq:kochkusche}
\Omega_k=\bigl\|(\mat L_k\hat{\vc x}-\vc x_{0,k})/\sigma_k\bigr\|^2,\qquad
r_k=p_k-\tr\bigl(\mat L_k\mat N^{-1}\mat L_k\T\bigr)/\sigma_k^2,\qquad
\sigma_k\leftarrow\sigma_k\sqrt{\Omega_k/r_k},
\end{equation}
so that $\lambda_k=\sigma_0^2/\sigma_k^2$ is the regularisation the data support. A group needs
several rows: with one row $r_k\le1$ and the component is not estimable. A prior on local
parameters is an observation stream of its own, and its weight is estimated as that stream's
level by \eqref{eq:simplified}.

\section{Skipped arcs}\label{sec:skips}

An arc without data, or whose covariance is not positive definite, is skipped in a pass and keeps
its previous factor. A pass in which more than a set share of a stream's arcs is skipped is refused
rather than estimated on a fraction of the data.

\section{What the data can resolve}

The variance of a component estimated from a redundancy $r_g$ is approximately
$2\theta_g^2/r_g$ under normality: a factor with $r_g=30$ is known to about $\pm26\,\%$. Few or
short arcs cannot carry their own factors; a spectral bin needs enough redundancy per pass, which
is what the bin threshold and the choice of the grid control. A spectral law describes stationary
noise; a periodic signal or a step in the residuals is a functional model error, which the
estimator would absorb into the noise at the price of pessimistic errors
(Chapter~\ref{ch:quality}).

% Chapter 4 -- robust estimation.
\chapter{Robust estimation}\label{ch:robust}

\section{Influence functions}

Robust estimation replaces the weight of an observation by $f(u)\,p$, where $u$ is its
standardised residual and $f$ the weight function of an M-estimator. The scale is the centred
median absolute deviation, $s=1.4826\,\operatorname{median}|v-\operatorname{median}v|$, and $u=v/s$.

\begin{center}
\begin{tabular}{@{}llll@{}}
\toprule
influence & $f(u)$ & default tuning & rejection point \\ \midrule
Huber & $1$ for $|u|\le c$, $c/|u|$ beyond & $c=1.345$ & none \\
IGG3 & $1$ on $[0,k_0]$, $\dfrac{k_0}{|u|}\Bigl(\dfrac{k_1-|u|}{k_1-k_0}\Bigr)^2$ on $(k_0,k_1]$, $0$ beyond & $k_0=1.5$, $k_1=3$ & $k_1$\\
bisquare & $\bigl(1-(u/c)^2\bigr)^2$ for $|u|<c$, $0$ beyond & $c=4.685$ & $c$ \\
\bottomrule
\end{tabular}
\end{center}

Huber is monotone and keeps the fraction $c/|u|$ of a gross error's weight; IGG3 and bisquare
redescend and reject it. Weights are floored at $10^{-12}$.

\section{Iteratively reweighted least squares on the innovations}\label{sec:irls}

Without variance-component estimation, IRLS weights the \emph{law-whitened} residuals (the
innovations of the noise process; Yang 1994). Each pass accumulates under the current factors,
solves, forms the whitened residuals with the previous factor removed, computes one MAD scale over
all data and the factors $f_i$, and applies them after the whitener: the whitened rows are scaled
by $\sqrt{f_i}$. The iteration stops when the change of $\hat{\vc x}$ falls below its tolerance
relative to the size of the estimate or its formal error. The returned $\hat\sigma_0$ and covariance
are those of the reweighted problem. For a correlated law a zero innovation weight removes the
whitening filter's response to the blunder, not the epoch itself.

\section{Robust weights inside the variance-component estimation}\label{sec:robustvce}

Inside the joint estimator the weights are statements about \emph{observations} (Koch's
convention), and they enter the covariance, not the whitened rows.

\paragraph{The statistic.} The raw residuals $\vc v=\vc l-\mat A\hat{\vc x}-\mat B\hat{\vc y}$ are
de-scaled by the declared per-epoch accuracy ($\mat D_e^{-1}\vc v_e$), not by the temporal law and not
by the level. For a scalar stream the statistic is the de-scaled residual. For an epoch of $d$ axes
it is $\chi$-calibrated,
\begin{equation}
z_e=\operatorname{sign}\Bigl(\sum_{\mathrm{axes}}v\Bigr)\,\Phi^{-1}\!\Bigl(1-\tfrac12\,Q\bigl(\tfrac d2,\tfrac{q_e}{2s^2}\bigr)\Bigr),\qquad
q_e=\sum_{\mathrm{axes}}v^2 ,
\end{equation}
with the upper regularised gamma function $Q$, so that a clean $d$-axis epoch is standardised like
a clean scalar residual. The scale is one MAD per stream, so the weights do not depend on the level
being estimated.

\paragraph{The nugget.} A factor $f_e$ enters as an additive variance on the diagonal of the
temporal law,
\begin{equation}\label{eq:nugget}
\mat\Sigma=\mat D\bigl((\mat R_t+\diag(\vc c))\otimes\mat I_d\bigr)\mat D\T,\qquad c_e=q_e\,(\mat R_t)_{ee},\quad q_e=\frac1{f_e}-1 .
\end{equation}
The marginal variance of the epoch becomes $(\mat R_t)_{ee}/f_e$, its covariances with the
neighbours are unchanged, and $f_e=0$ deletes the epoch exactly. The arc is whitened by the law's
factor followed by an exact rank-$k$ correction over the $k$ down-weighted epochs (Woodbury).

\paragraph{What uses which factors.} The solution uses all factors. The component update uses only
the rejected epochs: a partial factor would bias the level downwards, because the influence
function down-weights exactly the epochs with large residuals. The level and the arc factors are
divided by the truncation factor of the rejection point $k$,
\begin{equation}
\tau=\frac{F_{\chi^2_{d+2}}(q_k)}{F_{\chi^2_d}(q_k)},\qquad P(\chi^2_d>q_k)=P(|N(0,1)|>k),
\end{equation}
which is $1-2k\varphi(k)/(2\Phi(k)-1)=0.9733$ for $d=1$ and $k=3$, practically one for bisquare and
exactly one for Huber. After the loop the returned solution is the \emph{hard} one: rejected epochs
deleted, every partial factor set back to one, and $\hat\sigma_0^2$ and the covariance corrected for
the deleted rows and for $\tau$. A pass that rejects more than 45\,\% of the epochs has reached the
breakdown point of the estimator.

% Chapter 5 -- quality assessment.
\chapter{Quality assessment}\label{ch:quality}

All statistics of this chapter are formed in the whitened frame with a-priori $\sigma_0=1$; an
``observation'' is a whitened row, which for a correlated law is a linear combination of raw
epochs. The parameter set is the full one, $u=n_c+\sum_sn_{l,s}$, and the observations include the
prior rows, so $n-u$ is the total redundancy. The cofactor matrix is held as the blocks of
\eqref{eq:leverage}.

\section{The order of the verdict}

A result is judged in a fixed order, each level assuming the one before it holds:
\begin{enumerate}
\setcounter{enumi}{-1}
\item \textbf{Integrity.} The redundancy identity $\sum_ir_i=n-u$ holds, and the residual sum of
  the accumulation agrees with that of a direct pass over the residuals.
\item \textbf{Global test.} Is the model as a whole consistent with the declared noise?
\item \textbf{Localisation.} Which observations violate it?
\item \textbf{Reliability.} What could not have been detected, and how far could it move the
  estimates?
\item \textbf{Parameters.} What is actually known: precision, correlation, conditioning.
\end{enumerate}

\section{Global test}

$\hat\sigma_0^2/\sigma_0^2$ is tested against $\chi^2_{n-u}/(n-u)$ in the two-sided band
$[\chi^2_{0.005}/(n-u),\ \chi^2_{0.995}/(n-u)]$, with the 95\,\% confidence interval
$\bigl[\hat\sigma_0\sqrt{r/\chi^2_{0.975}(r)},\ \hat\sigma_0\sqrt{r/\chi^2_{0.025}(r)}\bigr]$. After
variance-component estimation the fixed point $\Omega_g=r_g$ makes the test pass by construction
for every estimated group.

\section{Redundancy numbers and the outlier tests}\label{sec:redundancy}

\paragraph{Redundancy numbers.} $r_i=1-\bar{\vc a}_i\T\mat Q_{xx}\bar{\vc a}_i$, evaluated exactly with
the local block \eqref{eq:leverage}. $r_i\to0$: the observation is fitted whatever its value (no
control); $r_i\to1$: it barely influences the solution.

\paragraph{Standardised residuals.} For an observation with $r_i>r_{\min}$,
\begin{equation}
w_i=\frac{\bar v_i}{\hat s\sqrt{r_i}},
\end{equation}
with $\hat s=\hat\sigma_0$ (Pope's test, $\sigma_0$ estimated), $\hat s=\sigma_0$ (Baarda's test,
$\sigma_0$ known) or a robust MAD scale; observations below $r_{\min}$ are uncontrolled and not
tested. Under Pope's test $w_i$ follows the $\tau$ distribution,
$\tau_r=t_{r-1}\sqrt r/\sqrt{r-1+t_{r-1}^2}$, bounded by $\sqrt r$, because $\hat\sigma_0$ contains the
tested residual. A robust scale does not: with the MAD, whose efficiency for a variance is 0.367,
$w_i$ follows Student's $t$ with $\nu=0.367\,r$ degrees of freedom. With $n$ tests the family-wise
level is controlled by Bonferroni, $\alpha'=\alpha/n$.

\paragraph{Iterative data snooping.} The observation with the largest $|w_i|$ above the critical
value is removed, the solution is updated, and the test is repeated with $\alpha$ divided by the
remaining number of observations. Removing whitened row $i$ with weight $p$ is a rank-one downdate:
with $\vc z=\vc a_c-\mat G\T\vc b$ and $\delta=1-p\,\vc b\T\mat N_{ll}^{-1}\vc b$,
\begin{equation}
\mat R_c\leftarrow\operatorname{cholupdate}\bigl(\mat R_c,\sqrt{p/\delta}\,\vc z,-\bigr),\qquad
\mat N_{ll}\leftarrow\mat N_{ll}-p\,\vc b\vc b\T,\qquad
\hat{\vc x}\leftarrow\hat{\vc x}+p\,\bar v_i\,\mat Q_{xx}\vc a_i ,
\end{equation}
and $\bar{\vc v}\T\bar{\vc v}$ decreases by $p\bar v_i^2/r_i$. Observations with low redundancy are not
rejected, and the number of rejections is capped; the final model is re-accumulated from the
surviving data.

\paragraph{Robust cross-check.} The raw scores $\rho_i=\bar v_i/\sqrt{r_i}$ give a robust scale
(the centred MAD). Their ratio to $\hat\sigma_0$, and the share $\hat f$ of scores beyond $2.5$
robust standard deviations (null value $2\Phi(-2.5)=0.0124$), gauge whether the least-squares scale
is inflated by contamination: a contaminated fit inflates $\hat\sigma_0$ and masks its own outliers.

\section{Reliability}\label{sec:reliability}

The minimal detectable bias of observation $i$ is
\begin{equation}
\nabla_{0,i}=\frac{\sigma_i\,\delta_0}{\sqrt{r_i}},\qquad
1-\Phi(k-\delta_0)+\Phi(-k-\delta_0)=1-\beta,\quad k=\Phi^{-1}(1-\alpha'/2),
\end{equation}
the two-sided power solved exactly at the Bonferroni-corrected level ($\delta_0=4.13$ for
$\alpha'=10^{-3}$, $\beta=0.2$). Its effect on the parameters (external reliability) is
$\nabla\hat{\vc x}=\mat Q_{xx}\vc a_i\,p_i\,\nabla_{0,i}$, reported relative to the parameters'
standard deviations. Observations with small $r_i$ --- short arcs, many local parameters --- are
where undetected errors hide.

\section{Whiteness and systematic errors}\label{sec:whiteness}

After variance-component estimation the whitened residuals of every stream should be white. The
Ljung--Box statistic pools the arcs of a stream: with lag sums $s_k=\sum\bar e_i\bar e_{i+k}$ and
pair counts $N_k=\sum_{\mathrm{arcs}}\max(m_a-k,0)$,
\begin{equation}
Q=N_0(N_0+2)\sum_{k=1}^{K}\frac{\rho_k^2}{N_k}\sim\chi^2_K,\qquad\rho_k=s_k/s_0,\quad K=20 .
\end{equation}
It detects correlation at short lags, periodic terms and excess low-frequency noise; a correlation
confined to a single long lag lies outside it. Structure in the residuals against time or another
covariate --- runs, binned means, periodogram peaks --- indicates a \emph{functional} model error.
Variance-component estimation cannot repair such an error; it absorbs it into the noise at the
price of pessimistic errors and down-weighted arcs.

\section{Parameter-level measures}

\paragraph{Correlation and conditioning.} Correlations follow from $\mat Q_{xx}$. The condition
number of $\mat N$ is bounded from below by power iterations on $\mat N$ and $\mat Q_{xx}$; it is a lower
bound.

\paragraph{Contribution of an observation group.} The share of group $g$ in the common solution is
its trace
\begin{equation}
t_g=\tr\bigl(\mat N^{-1}\mat N_g\bigr)=\sum_{i\in g}\vc z_i\T\mat N^{-1}\vc z_i ,
\end{equation}
the common term of the leverage \eqref{eq:leverage} summed over the group's rows; the $t_g$ of the
data sum to $n_c$. It is computed exactly in the redundancy sweep. The per-parameter contribution
$\bigl[\mat N^{-1}\mat N_g\bigr]_{jj}$, which sums to one over the groups and to $t_g$ over the
parameters, is estimated with random probes.

\paragraph{Nested models.} A full model with $u_F$ parameters against a reduced one with $u_R<u_F$ on
the same $n$ observations:
\begin{equation}
F=\frac{(\Omega_R-\Omega_F)/(u_F-u_R)}{\Omega_F/(n-u_F)}\sim F(u_F-u_R,\ n-u_F),
\end{equation}
where $u$ counts the common and the eliminated local parameters.

\paragraph{Confidence ellipses.} For a parameter pair, the ellipse of $\mat Q_{xx}$ restricted to the
pair, scaled by $k^2=2F_{1-\alpha}(2,r)\hat\sigma_0^2$ (or $\chi^2_{1-\alpha}(2)$ when $\sigma_0$ is known).

% References.
\chapter*{References}\addcontentsline{toc}{chapter}{References}
\markboth{References}{References}

\begin{description}[leftmargin=0pt,labelindent=0pt,itemsep=3pt]
\item Baarda, W. (1968). \emph{A testing procedure for use in geodetic networks.} Netherlands
  Geodetic Commission, Publications on Geodesy, 2(5).
\item Ellmer, M. (2018). \emph{Contributions to GRACE gravity field recovery: improvements in
  dynamic orbit integration, stochastic modelling of the antenna offset correction, and
  co-estimation of satellite orientations.} Monographic Series TU Graz, Graz University of
  Technology.
\item Förstner, W. (1979). Ein Verfahren zur Schätzung von Varianz- und Kovarianzkomponenten.
  \emph{Allgemeine Vermessungs-Nachrichten} 86, 446--453.
\item Hutchinson, M. F. (1990). A stochastic estimator of the trace of the influence matrix for
  Laplacian smoothing splines. \emph{Communications in Statistics -- Simulation and Computation}
  19(2), 433--450.
\item Koch, K.-R. (1999). \emph{Parameter estimation and hypothesis testing in linear models.}
  2nd ed., Springer.
\item Koch, K.-R., Kusche, J. (2002). Regularization of geopotential determination from satellite
  data by variance components. \emph{Journal of Geodesy} 76, 259--268.
\item Kvas, A. (2019). \emph{Gravity field recovery from GRACE-type observations: daily solutions
  and process noise.} PhD thesis, Graz University of Technology.
\item Ljung, G. M., Box, G. E. P. (1978). On a measure of lack of fit in time series models.
  \emph{Biometrika} 65(2), 297--303.
\item Pope, A. J. (1976). \emph{The statistics of residuals and the detection of outliers.} NOAA
  Technical Report NOS 65 NGS 1.
\item Yang, Y. (1994). Robust estimation for dependent observations. \emph{Manuscripta Geodaetica}
  19, 10--17.
\end{description}


\end{document}
